\ifdefined\gkver\else\def\gkver{student}\fi
\documentclass[\gkver]{gaokaozhenti}
\newcommand{\ccnum}[1]{{\fontspec{Noto Serif CJK SC}#1}}
\begin{document}

\chapter{2014年山东理综}
%% paperType: 理综物理部分

\begin{enumerate}
\item
%% number: 14
%% paperName: 2014年山东理综第 14 题 6 分
%% typeId: 1
%% type: 单选题
%% chapter: 力物体的平衡
%% point: 三力平衡
%% method: 无
%% score: 6
%% degree: 650
%% duplicateId: 0
%% body:
如图，用两根等长轻绳将木板悬挂在竖直木桩上等高的两点，制成一简易秋千。某次维修时将两轻绳各剪去一小段，但仍保持等长且悬挂点不变。木板静止时，$F_{1}$ 表示木板所受合力的大小，$F_{2}$ 表示单根轻绳对木板拉力的大小，则维修后\xzanswer{A}

\onepicture[w=3.0cm]{figs/14.png}

\fourchoices[answer=A]
{$F_{1}$ 不变，$F_{2}$ 变大}
{$F_{1}$ 不变，$F_{2}$ 变小}
{$F_{1}$ 变大，$F_{2}$ 变大}
{$F_{1}$ 变小，$F_{2}$ 变小}

%% answer: A
%% memo:
\memoanswer{
由于木板始终处于静止状态，因此木板所受合力为零，故选项 C、D 错误；对木板进行受力分析如图所示，由平衡条件得：$2F_{2}\cos\theta=G$，当轻绳被剪短后，$\theta$ 增大，$\cos\theta$ 减小，则 $F_{2}$ 增大，故选项 A 正确，B 错误。

\onepicture[w=5.0cm]{figs/14_an.png}
}

\item
%% number: 15
%% paperName: 2014年山东理综第 15 题 6 分
%% typeId: 2
%% type: 多选题
%% chapter: 运动和力
%% point: 牛顿第二定律
%% method: 无
%% score: 6
%% degree: 650
%% duplicateId: 0
%% body:
一质点在外力作用下做直线运动，其速度 $v$ 随时间 $t$ 变化的图像如图。在图中标出的时刻中，质点所受合外力的方向与速度方向相同的有\xzanswer{AC}

\onepicture[w=5.0cm]{figs/15.png}

\fourchoices[answer=AC]
{$t_{1}$}
{$t_{2}$}
{$t_{3}$}
{$t_{4}$}

%% answer: AC
%% memo:
\memoanswer{
当合外力方向与速度方向相同时，质点做加速运动。由 $v-t$ 图像可知，质点在 $t_{1}$、$t_{3}$ 时刻做加速运动，在 $t_{2}$、$t_{4}$ 时刻做减速运动。故选项 A、C 正确，选项 B、D 错误。
}

\item
%% number: 16
%% paperName: 2014年山东理综第 16 题 6 分
%% typeId: 2
%% type: 多选题
%% chapter: 电磁感应
%% point: 电磁感应受力分析
%% method: 无
%% score: 6
%% degree: 650
%% duplicateId: 0
%% body:
如图，一端接有定值电阻的平行金属轨道固定在水平面内，通有恒定电流的长直绝缘导线垂直并紧靠轨道固定，导体棒与轨道垂直且接触良好。在向右匀速通过 M、N 两区的过程中，导体棒所受安培力分别用 $F_{M}$、$F_{N}$ 表示。不计轨道电阻。以下叙述正确的是\xzanswer{BCD}

\onepicture[w=5.5cm]{figs/16.png}

\fourchoices[answer=BCD]
{$F_{M}$ 向右}
{$F_{N}$ 向左}
{$F_{M}$ 逐渐增大}
{$F_{M}$ 逐渐减小}

%% answer: BCD
%% memo:
\memoanswer{
本题原卷及材料中未提供详解，待补充。
}

\item
%% number: 17
%% paperName: 2014年山东理综第 17 题 6 分
%% typeId: 2
%% type: 多选题
%% chapter: 交流电
%% point: 变压器
%% method: 无
%% score: 6
%% degree: 650
%% duplicateId: 0
%% body:
如图，将额定电压为 $60\UV$ 的用电器，通过一理想变压器接在正弦交变电源上。闭合开关 S 后，用电器正常工作，交流电压表和交流电流表（均为理想电表）的示数分别为 $220\UV$ 和 $2.2\UA$。以下判断正确的是\xzanswer{BD}

\onepicture[w=5.5cm]{figs/17.png}

\fourchoices[answer=BD]
{变压器输入功率为 $484\UW$}
{通过原线圈的电流的有效值为 $6.6\UA$}
{通过副线圈的电流的最大值为 $2.2\UA$}
{变压器原、副线圈匝数比 $n_{1}:n_{2}=11:3$}

%% answer: BD
%% memo:
\memoanswer{
本题原卷及材料中未提供详解，待补充。
}

\item
%% number: 18
%% paperName: 2014年山东理综第 18 题 6 分
%% typeId: 1
%% type: 单选题
%% chapter: 电场
%% point: 带电粒子的偏转
%% method: 无
%% score: 6
%% degree: 450
%% duplicateId: 0
%% body:
如图，场强大小为 $E$、方向竖直向下的匀强电场中有一矩形区域 abcd，水平边 ab 长为 $s$，竖直边 ad 长为 $h$。质量均为 $m$、带电量分别为 $+q$ 和 $-q$ 的两粒子，由 a、c 两点先后沿 ab 和 cd 方向以速率 $v_{0}$ 进入矩形区（两粒子不同时出现在电场中）。不计重力。若两粒子轨迹恰好相切，则 $v_{0}$ 等于\xzanswer{B}

\onepicture[w=4.5cm]{figs/18.png}

\fourchoices[answer=B]
{$\frac{s}{2}\sqrt{\frac{2qE}{mh}}$}
{$\frac{s}{2}\sqrt{\frac{qE}{mh}}$}
{$\frac{s}{4}\sqrt{\frac{2qE}{mh}}$}
{$\frac{s}{4}\sqrt{\frac{qE}{mh}}$}

%% answer: B
%% memo:
\memoanswer{
本题原卷及材料中未提供详解，待补充。
}

\item
%% number: 19
%% paperName: 2014年山东理综第 19 题 6 分
%% typeId: 1
%% type: 单选题
%% chapter: 电场
%% point: 带电粒子的直线运动
%% method: 无
%% score: 6
%% degree: 450
%% duplicateId: 0
%% body:
如图，半径为 $R$ 的均匀带正电的薄球壳，其上有一小孔 A。已知壳内的场强处处为零；壳外空间的电场，与将球壳上的全部电荷集中于球心 O 时在壳外产生的电场一样。一带正电的试探电荷（不计重力）从球心以初动能 $E_{k0}$ 沿 OA 方向射出。下列关于试探电荷的动能 $E_{k}$ 与离开球心的距离 $r$ 的关系，可能正确的是\xzanswer{A}

\onepicture[w=3.0cm]{figs/19a.png}

\fourchoices[answer=A, ispicture=true, h=2.0cm]
{figs/19b.png}
{figs/19c.png}
{figs/19d.png}
{figs/19e.png}

%% answer: A
%% memo:
\memoanswer{
本题原卷及材料中未提供详解，待补充。
}

\item
%% number: 20
%% paperName: 2014年山东理综第 20 题 6 分
%% typeId: 1
%% type: 单选题
%% chapter: 曲线运动
%% point: 万有引力定律
%% method: 无
%% score: 6
%% degree: 450
%% duplicateId: 0
%% body:
2013 年我国相继完成“神十”与“天宫”对接、“嫦娥”携“玉兔”落月两大航天工程。某航天爱好者提出“玉兔”回家的设想：如图，将携带“玉兔”的返回系统由月球表面发射到 $h$ 高度的轨道上，与在该轨道绕月球做圆周运动的飞船对接，然后由飞船送“玉兔”返回地球。设“玉兔”质量为 $m$，月球半径为 $R$，月面的重力加速度为 $g_{\text{月}}$。以月面为零势能面，“玉兔”在 $h$ 高度的引力势能可表示为 $E_{p}=\frac{GMmh}{R(R+h)}$，其中 $G$ 为引力常量，$M$ 为月球质量。若忽略月球的自转，从开始发射到对接完成需要对“玉兔”做的功为\xzanswer{D}

\onepicture[w=4.0cm]{figs/20.png}

\fourchoices[answer=D]
{$\frac{mg_{\text{月}}R}{R+h}(h+2R)$}
{$\frac{mg_{\text{月}}R}{R+h}(h+\sqrt{2}R)$}
{$\frac{mg_{\text{月}}R}{R+h}(h+\frac{\sqrt{2}}{2}R)$}
{$\frac{mg_{\text{月}}R}{R+h}(h+\frac{1}{2}R)$}

%% answer: D
%% memo:
\memoanswer{
本题原卷及材料中未提供详解，待补充。
}

\item
%% number: 21
%% paperName: 2014年山东理综第 21 题 8 分
%% typeId: 4
%% type: 实验题
%% chapter: 功和能
%% point: 动能定理
%% method: 整体与隔离
%% score: 8
%% degree: 650
%% duplicateId: 0
%% body:
某实验小组利用弹簧秤和刻度尺，测量滑块在木板上运动的最大速度。实验步骤：
\begin{steps}
	\item
	用弹簧秤测量橡皮泥和滑块的总重力，记作 $G$；
	\item
	将装有橡皮泥的滑块放在水平木板上，通过水平细绳和固定弹簧秤相连，如图甲所示。在 A 端向右拉动木板，等弹簧秤示数稳定后，将计数记作 $F$；
	\item
	改变滑块上橡皮泥的质量，重复步骤①②；实验数据如下表所示：

\begin{center}
\begin{tabular}{|c|c|c|c|c|c|c|}
\hline
$G/\UN$ & 1.50 & 2.00 & 2.50 & 3.00 & 3.50 & 4.00 \\ \hline
$F/\UN$ & 0.59 & 0.83 & 0.99 & 1.22 & 1.37 & 1.61 \\ \hline
\end{tabular}
\end{center}

	\item
	如图乙所示，将木板固定在水平桌面上，滑块置于木板上左端 C 处，细绳跨过定滑轮分别与滑块和重物 P 连接，保持滑块静止，测量重物 P 离地面的高度 $h$；
	\item
	滑块由静止释放后开始运动并最终停在木板上的 D 点（未与滑轮碰撞），测量 C、D 间的距离 $s$。
\end{steps}

完成下列作图和填空：
\begin{enumerate}
	\item
	根据表中数据在给定的坐标纸（见答题卡）上作出 $F-G$ 图线。
	\item
	由图线求得滑块和木板间的动摩擦因数 $\mu=$\_\_\_\_\_\_（保留 2 位有效数字）
	\item
	滑块最大速度的大小 $v=$\_\_\_\_\_\_（用 $h$、$s$、$\mu$ 和重力加速度 $g$ 表示。）
\end{enumerate}

\twopicture[numstyle=chinese, labelA=2014山东理综21a, labelB=2014山东理综21b, wA=6.0cm, wB=5.5cm, align=right]{figs/21a.png}{figs/21b.png}

\onepicture[w=4.5cm, align=right]{figs/21c.png}

%% answer:
\jdanswer*{
\begin{enumerate}
	\item
	如图所示

\onepicture[w=4.5cm]{figs/21_an.png}

	\item
	0.40（0.38、0.39、0.41、0.42 均正确）
	\item
	$2\mu g(s-h)$
\end{enumerate}
}
%% memo:
\memoanswer{
本题原卷及材料中未提供详解，待补充。
}

\item
%% number: 22
%% paperName: 2014年山东理综第 22 题 10 分
%% typeId: 4
%% type: 实验题
%% chapter: 电路
%% point: 电路实验
%% method: 无
%% score: 10
%% degree: 450
%% duplicateId: 0
%% body:
实验室购买了一捆标称长度为 $100\Um$ 的铜导线，某同学想通过实验测定其实际长度。该同学首先测得导线横截面积为 $1.0\Ummq$，查得铜的电阻率为 $1.7\times10^{-8}\UOm$，再利用图甲所示电路测出铜导线的电阻 $R_{x}$，从而确定导线的实际长度。

可供使用的器材有：

电流表：量程 $0.6\UA$，内阻约为 $0.2\UO$；

电压表：量程 $3\UV$，内阻约 $9\UkO$；

滑动变阻器 $R_{1}$：最大阻值 $5\UO$；

滑动变阻器 $R_{2}$：最大阻值 $20\UO$；

定值电阻：$R_{0}=3\UO$；

电源：电动势 $6\UV$，内阻可不计；

开关、导线若干。

回答下列问题：
\begin{enumerate}
	\item
	实验中滑动变阻器应选\_\_\_\_\_\_（填“$R_{1}$”或“$R_{2}$”），闭合开关 S 前应将滑片移至\_\_\_\_\_\_端（填“a”或“b”）。
	\item
	在实物图中，已正确连接了部分导线，请根据图甲电路完成剩余部分的连接。
	\item
	调节滑动变阻器，当电流表的读数为 $0.50\UA$ 时，电压表示数如图乙所示，读数为\_\_\_\_\_\_$\UV$。
	\item
	导线实际长度为\_\_\_\_\_\_$\Um$（保留 2 位有效数字）。
\end{enumerate}

\twopicture[numstyle=chinese, labelA=2014山东理综22a, labelB=2014山东理综22b, wA=6.0cm, wB=5.0cm, align=right]{figs/22a.png}{figs/22b.png}

\onepicture[w=5.0cm, align=right]{figs/22c.png}

%% answer:
\jdanswer*{
\begin{enumerate}
	\item
	$R_{2}$；a
	\item
	如图所示

\onepicture[w=5.0cm]{figs/22_an.png}

	\item
	2.30（2.29、2.31 均正确）
	\item
	94（93、95 均正确）
\end{enumerate}
}
%% memo:
\memoanswer{
本题原卷及材料中未提供详解，待补充。
}

\item
%% number: 23
%% paperName: 2014年山东理综第 23 题 18 分
%% typeId: 6
%% type: 计算题
%% chapter: 运动和力
%% point: 牛顿定律的应用
%% method: 无
%% score: 18
%% degree: 650
%% duplicateId: 0
%% body:
研究表明，一般人的刹车反应时间（即图甲中“反应过程”所用时间）$t_{0}=0.4\Us$，但饮酒会导致反应时间延长。在某次试验中，志愿者少量饮酒后驾车以 $v_{0}=72\Ukmh$ 的速度在试验场的水平路面上匀速行驶，从发现情况到汽车停止，行驶距离 $L=39\Um$。减速过程中汽车位移 $s$ 与速度 $v$ 的关系曲线如图乙所示，此过程可视为匀变速直线运动。取重力加速度的大小 $g=10\Umsq$。求：
\begin{enumerate}
	\item
	减速过程汽车加速度的大小及所用时间；
	\item
	饮酒使志愿者的反应时间比一般人增加了多少？
	\item
	减速过程汽车对志愿者作用力的大小与志愿者重力大小的比值。
\end{enumerate}

\twopicture[numstyle=chinese, labelA=2014山东理综23a, labelB=2014山东理综23b, wA=6.5cm, wB=4.5cm, align=right]{figs/23a.png}{figs/23b.png}

%% answer:
\jdanswer{
\begin{enumerate}
	\item
	$a=8\Umsq$，$t=2.5\Us$
	\item
	$\Delta t=0.3\Us$
	\item
	$\frac{\sqrt{41}}{5}$
\end{enumerate}
}
%% memo:
\memoanswer{
（1）设减速过程中汽车加速度的大小为 $a$，所用时间为 $t$，由题可得初速度 $v_{0}=20\Ums$，末速度 $v_{t}=0$，位移 $s=25\Um$，由运动学公式得

$v_{0}^{2}=2as$ ①

$t=\frac{v_{0}}{a}$ ②

联立①②式，代入数据得

$a=8\Umsq$ ③

$t=2.5\Us$ ④

（2）设志愿者反应时间为 $t'$，反应时间的增加量为 $\Delta t$，由运动学公式得

$L=v_{0}t'+s$ ⑤

$\Delta t=t'-t_{0}$ ⑥

联立⑤⑥式，代入数据得

$\Delta t=0.3\Us$ ⑦

（3）设志愿者所受合外力的大小为 $F$，汽车对志愿者作用力的大小为 $F_{0}$，志愿者质量为 $m$，由牛顿第二定律得

$F=ma$ ⑧

由平行四边形定则得

$F_{0}^{2}=F^{2}+(mg)^{2}$ ⑨

联立③⑧⑨式，代入数据得

$\frac{F_{0}}{mg}=\frac{\sqrt{41}}{5}$ ⑩
}

\item
%% number: 24
%% paperName: 2014年山东理综第 24 题 20 分
%% typeId: 6
%% type: 计算题
%% chapter: 磁场
%% point: 带电粒子在复合场中的运动
%% method: 无
%% score: 20
%% degree: 250
%% duplicateId: 0
%% body:
如图甲所示，间距为 $d$ 且垂直于纸面的两平行板 P、Q 间存在匀强磁场。取垂直于纸面向里为磁场的正方向，磁感应强度随时间的变化规律如图乙所示。$t=0$ 时刻，一质量为 $m$、带电量为 $+q$ 的粒子（不计重力），以初速度 $v_{0}$ 由 Q 板左端靠近板面的位置，沿垂直于磁场且平行于板面的方向射入磁场区。当 $B_{0}$ 和 $T_{B}$ 取某些特定值时，可使 $t=0$ 时刻入射的粒子经 $\Delta t$ 时间恰能垂直打在 P 板上（不考虑粒子反弹）。上述 $m$、$q$、$d$、$v_{0}$ 为已知量。
\begin{enumerate}
	\item
	若 $\Delta t=\frac{1}{2}T_{B}$，求 $B_{0}$；
	\item
	若 $\Delta t=\frac{3}{2}T_{B}$，求粒子在磁场中运动时加速度的大小；
	\item
	若 $B_{0}=\frac{4mv_{0}}{qd}$，为使粒子仍能垂直打在 P 板上，求 $T_{B}$。
\end{enumerate}

\twopicture[numstyle=chinese, labelA=2014山东理综24a, labelB=2014山东理综24b, wA=4.5cm, wB=5.5cm, align=right]{figs/24a.png}{figs/24b.png}

%% answer:
\jdanswer{
\begin{enumerate}
	\item
	$B_{0}=\frac{mv_{0}}{qd}$
	\item
	$a=\frac{3v_{0}^{2}}{d}$
	\item
	$T_{B}=\frac{\pi d}{3v_{0}}$ 或 $T_{B}=\left(\frac{\pi}{2}+\arcsin\frac{1}{4}\right)\frac{d}{2v_{0}}$
\end{enumerate}
}
%% memo:
\memoanswer{
（1）设粒子做圆周运动的半径为 $R_{1}$，由牛顿第二定律得

$qv_{0}B_{0}=\frac{mv_{0}^{2}}{R_{1}}$ ①

据题意由几何关系得 $R_{1}=d$ ②

联立①②式得 $B_{0}=\frac{mv_{0}}{qd}$ ③

（2）设粒子做圆周运动的半径为 $R_{2}$，加速度大小为 $a$，由圆周运动公式得

$a=\frac{v_{0}^{2}}{R_{2}}$ ④

据题意由几何关系得 $3R_{2}=d$ ⑤

联立④⑤式得 $a=\frac{3v_{0}^{2}}{d}$ ⑥

（3）设粒子做圆周运动的半径为 $R$，周期为 $T$，由圆周运动公式得

$T=\frac{2\pi R}{v_{0}}$ ⑦

由牛顿第二定律得 $qv_{0}B_{0}=\frac{mv_{0}^{2}}{R}$ ⑧

由题意知 $B_{0}=\frac{4mv_{0}}{qd}$，代入⑧式得

$d=4R$ ⑨

粒子运动轨迹如图所示，$O_{1}$、$O_{2}$ 为圆心，$O_{1}O_{2}$ 连线与水平方向夹角为 $\theta$。在第 $T_{B}$ 内，只有 A、B 两个位置才有可能垂直击中 P 板，且均要求 $0<\theta<\frac{\pi}{2}$，由题意可知

$\frac{\frac{\pi}{2}+\theta}{2\pi}T=\frac{T_{B}}{2}$ ⑩

\onepicture[w=4.5cm]{figs/24_an.png}

设经历整个完整 $T_{B}$ 的个数为 $n$（$n=0$、1、2、3……）

若在 A 点击中 P 板，据题意由几何关系得

$R+2(R+R\sin\theta)n=d$ \ccnum{⑪}

当 $n=0$ 时，无解 \ccnum{⑫}

当 $n=1$ 时，联立⑨\ccnum{⑪}式得

$\theta=\frac{\pi}{6}$（或 $\sin\theta=\frac{1}{2}$） \ccnum{⑬}

联立⑦⑨⑩\ccnum{⑬}式得 $T_{B}=\frac{\pi d}{3v_{0}}$ \ccnum{⑭}

当 $n\geq2$ 时，不满足 $0<\theta<\frac{\pi}{2}$ 的要求 \ccnum{⑮}

若在 B 点击中 P 板，据题意由几何关系得

$R+2R\sin\theta+2(R+R\sin\theta)n=d$ \ccnum{⑯}

当 $n=0$ 时，无解 \ccnum{⑰}

当 $n=1$ 时，联立⑨\ccnum{⑯}式得

$\theta=\arcsin\frac{1}{4}$（或 $\sin\theta=\frac{1}{4}$） \ccnum{⑱}

联立⑦⑨⑩\ccnum{⑱}式得 $T_{B}=\left(\frac{\pi}{2}+\arcsin\frac{1}{4}\right)\frac{d}{2v_{0}}$ \ccnum{⑲}

当 $n\geq2$ 时，不满足 $0<\theta<\frac{\pi}{2}$ 的要求 \ccnum{⑳}
}

\item
%% number: 37
%% paperName: 2014年山东理综第 37 题 7 分
%% typeId: 6
%% type: 计算题
%% chapter: 气体性质
%% point: 玻意耳定律
%% method: 无
%% score: 7
%% degree: 650
%% duplicateId: 0
%% body:
\textbf{（1）}（5 分）如图，内壁光滑、导热良好的气缸中用活塞封闭有一定质量的理想气体。当环境温度升高时，缸内气体\xzanswer{AB}。（双选，填正确答案标号）

\fourchoices[answer=AB]
{内能增加}
{对外做功}
{压强增大}
{分子间的引力和斥力都增大}

\onepicture[w=2.6cm]{figs/37a.png}

\textbf{（2）}（7 分）一种水下重物打捞方法的工作原理如图所示。将一质量 $M=3\times10^{3}\Ukg$、体积 $V_{0}=0.5\Umc$ 的重物捆绑在开口向下的浮筒上，向浮筒内充入一定量的气体。开始时筒内液面到水面的距离 $h_{1}=40\Um$，筒内气体体积 $V_{1}=1\Umc$，在拉力作用下浮筒缓慢上升，当筒内液面到水面的距离为 $h_{2}$ 时，拉力减为零，此时气体体积为 $V_{2}$，随后浮筒和重物自动上浮。求 $V_{2}$ 和 $h_{2}$。

已知大气压强 $p_{0}=1\times10^{5}\UPa$，水的密度 $\rho=1\times10^{3}\Ukgmc$，重力加速度的大小 $g=10\Umsq$，不计水温变化，筒内气体质量不变且可看作理想气体，浮筒质量和筒壁厚度可忽略。

\onepicture[w=4.0cm, align=right]{figs/37b.png}

%% answer:
\jdanswer{
$V_{2}=2.5\Umc$，$h_{2}=10\Um$
}
%% memo:
\memoanswer{
（2）当 $F=0$ 时，由平衡条件得

$Mg=\rho g(V_{0}+V_{2})$ ①

代入数据得

$V_{2}=2.5\Umc$ ②

设筒内气体初态、末态的压强分别为 $P_{1}$、$P_{2}$，由题意得

$P_{1}=P_{0}+\rho gh_{1}$ ③

$P_{2}=P_{0}+\rho gh_{2}$ ④

在此过程中筒内气体温度和质量不变，由玻意耳定律得

$P_{1}V_{1}=P_{2}V_{2}$ ⑤

联立②③④⑤式，代入数据得

$h_{2}=10\Um$ ⑥
}

\item
%% number: 38
%% paperName: 2014年山东理综第 38 题 7 分
%% typeId: 6
%% type: 计算题
%% chapter: 光学
%% point: 全反射
%% method: 无
%% score: 7
%% degree: 650
%% duplicateId: 0
%% body:
\textbf{（1）}（5 分）一列简谐横波沿直线传播。以波源 O 由平衡位置开始振动为计时零点，质点 A 的振动图像如图所示，已知 O、A 的平衡位置相距 $0.9\Um$。以下判断正确的是\xzanswer{AB}。（双选，填正确答案标号）

\fourchoices[answer=AB]
{波长为 $1.2\Um$}
{波源起振方向沿 $y$ 轴正方向}
{波速大小为 $0.4\Ums$}
{质点 A 的动能在 $t=4\Us$ 时最大}

\onepicture[w=5.0cm]{figs/38a.png}

\textbf{（2）}（7 分）如图，三角形 ABC 为某透明介质的横截面，O 为 BC 边的中点，位于截面所在平面内的一束光线自 O 以角 $i$ 入射，第一次到达 AB 边恰好发生全反射。已知 $\theta=15^{\circ}$，BC 边长为 $2L$，该介质的折射率为 $\sqrt{2}$。求：
\begin{enumerate}
	\item
	入射角 $i$；
	\item
	从入射到发生第一次全反射所用的时间（设光在真空中的速度为 $c$，可能用到 $\sin75^{\circ}=\frac{\sqrt{6}+\sqrt{2}}{4}$ 或 $\tan15^{\circ}=2-\sqrt{3}$）。
\end{enumerate}

\onepicture[w=3.5cm, align=right]{\includesvg{figs/38b}}

%% answer:
\jdanswer{
\begin{enumerate}
	\item
	$i=45^{\circ}$
	\item
	$\frac{(\sqrt{6}+\sqrt{2})L}{2c}$
\end{enumerate}
}
%% memo:
\memoanswer{
（2）(i) 根据全反射定律可知，光线在 AB 面上 P 点的入射角等于临界角 $C$，由折射定律得

$\sin C=\frac{1}{n}$ ①

代入数据得 $C=45^{\circ}$ ②

设光线在 BC 面上的折射角为 $r$，由几何关系得

$r=30^{\circ}$ ③

由折射定律得 $n=\frac{\sin i}{\sin r}$ ④

联立③④式，代入数据得 $i=45^{\circ}$ ⑤

\onepicture[w=4.0cm]{figs/38_an.png}

(ii) 在 $\triangle OPB$ 中，根据正弦定理得

$\frac{\overline{OP}}{\sin75^{\circ}}=\frac{L}{\sin45^{\circ}}$ ⑥

设所用时间为 $t$，光线在介质中的速度为 $v$，得

$\overline{OP}=vt$ ⑦

$v=\frac{c}{n}$ ⑧

联立⑥⑦⑧式，代入数据得

$t=\frac{(\sqrt{6}+\sqrt{2})L}{2c}$ ⑨
}

\item
%% number: 39
%% paperName: 2014年山东理综第 39 题 7 分
%% typeId: 6
%% type: 计算题
%% chapter: 运动和力
%% point: 动量守恒定律
%% method: 无
%% score: 7
%% degree: 650
%% duplicateId: 0
%% body:
\textbf{（1）}（5 分）氢原子能级如图，当氢原子从 $n=3$ 跃迁到 $n=2$ 的能级时，辐射光的波长为 $656\Unm$。以下判断正确的是\xzanswer{CD}。（双选，填正确答案标号）

\onepicture[w=3.5cm]{figs/39a.png}

\fourchoices[answer=CD]
{氢原子从 $n=2$ 跃迁到 $n=1$ 的能级时，辐射光的波长大于 $656\Unm$}
{用波长为 $325\Unm$ 的光照射，可使氢原子从 $n=1$ 跃迁到 $n=2$ 的能级}
{一群处于 $n=3$ 能级上的氢原子向低能级跃迁时最多产生 3 种谱线}
{用波长为 $633\Unm$ 的光照射，不能使氢原子从 $n=2$ 跃迁到 $n=3$ 的能级}

\textbf{（2）}（7 分）如图，光滑水平直轨道上两滑块 A、B 用橡皮筋连接，A 的质量为 $m$。开始时橡皮筋松弛，B 静止，给 A 向左的初速度 $v_{0}$。一段时间后，B 与 A 同向运动发生碰撞并粘在一起。碰撞后的共同速度是碰撞前瞬间 A 的速度的两倍，也是碰撞前瞬间 B 的速度的一半。求：
\begin{enumerate}
	\item
	B 的质量；
	\item
	碰撞过程中 A、B 系统机械能的损失。
\end{enumerate}

\onepicture[w=6.0cm, align=right]{figs/39b.png}

%% answer:
\jdanswer{
\begin{enumerate}
	\item
	$m_{B}=\frac{m}{2}$
	\item
	$\Delta E=\frac{1}{6}mv_{0}^{2}$
\end{enumerate}
}
%% memo:
\memoanswer{
（2）(i) 以初速度 $v_{0}$ 的方向为正方向，设 B 的质量为 $m_{B}$，A、B 碰撞后的共同速度为 $v$，由题意知：碰撞前瞬间 A 的速度为 $\frac{v}{2}$，碰撞前瞬间 B 的速度为 $2v$，由动量守恒定律得

$m\frac{v}{2}+2m_{B}v=(m+m_{B})v$ ①

由①式得 $m_{B}=\frac{m}{2}$ ②

(ii) 从开始到碰后的全过程，由动量守恒定律得

$mv_{0}=(m+m_{B})v$ ③

设碰撞过程 A、B 系统机械能的损失为 $\Delta E$，则

$\Delta E=\frac{1}{2}m\left(\frac{v}{2}\right)^{2}+\frac{1}{2}m_{B}(2v)^{2}-\frac{1}{2}(m+m_{B})v^{2}$ ④

联立②③④式得 $\Delta E=\frac{1}{6}mv_{0}^{2}$ ⑤
}

\end{enumerate}

\end{document}
